\begin{banner}
REMEDIËRING WISKUNDE \quad·\quad 6 OPVOEDING EN BEGELEIDING
\end{banner}

\begin{titel}
Statistiek
\end{titel}

\begin{ondertitel}
Van data naar besluit: jouw route in 7 levels
\end{ondertitel}

Als begeleider, opvoeder of zorgverlener krijg je voortdurend cijfers voor je neus: tevredenheidsenquêtes bij ouders, wachttijden, de resultaten van een nieuwe aanpak, een krantenkop met ‘uit onderzoek blijkt…’. Met deze bundel leer je die cijfers \textbf{beschrijven}, \textbf{kritisch bekijken} en er \textbf{correcte besluiten} uit trekken.

Je speelt de bundel door zoals een spel: je start in level 0 en werkt je level per level op. Aan het einde van elk level test je jezelf met een \textbf{CHECK}. Haal je de score? Dan ga je door naar het volgende level. Lukt het nog niet? Dan herneem je de theorie en de oefeningen en probeer je opnieuw. Op het einde wacht de \textbf{eindtoets}.

\textbf{Naam:} \_\_\_\_\_\_\_\_\_\_\_\_\_\_\_\_\_\_\_\_\_\_\_\_\_\_\_\_\_\_ \quad \textbf{Klas:} \_\_\_\_\_\_\_\_\_\_\_\_\_\_ \quad \textbf{Datum:} \_\_\_\_\_\_\_\_\_\_\_\_\_\_

\subsection*{Jouw route}

\breedtes{12/50/16/12}
\begin{tabular}{llll}
\hline
Level & Onderwerp & Leerplandoel & Behaald? \\
\hline
0 & Startklaar: de basisbegrippen & springplank & ☐ \\
1 & Centrummaten: gemiddelde, mediaan, modus & springplank & ☐ \\
2 & Spreidingsmaten: R, ΔQ, variantie, standaardafwijking & springplank & ☐ \\
3 & Steekproef en representativiteit & LPD 11 & ☐ \\
4 & Samenhang of oorzaak? Causaliteit en onderzoeksopzet & LPD 12 & ☐ \\
5 & De normale verdeling als model & LPD 13 & ☐ \\
6 & Kansen berekenen bij de normale verdeling & LPD 14 & ☐ \\
7 & Eindmissie: problemen oplossen, ICT en communiceren & LPD 1 – 4 & ☐ \\
\textbf{Eind} & \textbf{Eindtoets} & alle doelen & ☐ \\
\hline
\end{tabular}


\pagina

\section{Zo werk je met deze bundel}

\begin{missie}
Missie: weten wat je moet doen, en hoe je het doet.
\end{missie}

\begin{theorie}
\textbf{THEORIE} \quad Hier staat de leerstof, steeds met een voorbeeld erbij. Lees ze rustig en werk mee met potlood en rekenmachine.
\end{theorie}

\begin{voorbeeld}
\textbf{VOORBEELD} \quad Een uitgewerkt voorbeeld. Kijk goed hoe elke stap wordt opgebouwd; daarna doe jij hetzelfde.
\end{voorbeeld}

\begin{letop}
\textbf{LET OP} \quad Hier gaan leerlingen vaak de mist in. Lees deze kaders extra goed.
\end{letop}

\begin{ict}
\textbf{ICT-TIP} \quad Zo laat je je rekenmachine, GeoGebra of Excel het rekenwerk doen.
\end{ict}

\begin{doel}
\textbf{DOEL} \quad Dit moet je kunnen na het level. Kruis het vakje aan wanneer het lukt.
\end{doel}

\subsection*{De oefeningen}

Elke oefening heeft een titel en een moeilijkheidsgraad: \textbf{★☆☆} (opwarmer, met stap-voor-stap hulp), \textbf{★★☆} (gewone oefening) en \textbf{★★★} (uitdaging). Je vindt ook steeds deze soorten opdrachten terug:

\breedtes{22/70}
\begin{tabular}{ll}
\hline
Soort & Wat moet je doen? \\
\hline
\textbf{Aan de slag} & Een oefening waarbij je de theorie toepast op getallen. \\
\textbf{Context} & Een situatie uit de opvoedings-, zorg- of welzijnswereld. Je moet zelf nagaan welk begrip je nodig hebt. \\
\textbf{Invullen} & Een tekst met ontbrekende woorden. Kies uit de woordenlijst. \\
\textbf{Koppel} & Je koppelt een begrip aan een situatie, een grafiek of een omschrijving. \\
\textbf{Foutjager} & Iemand heeft een fout gemaakt. Zoek ze, leg uit en verbeter. \\
\textbf{CHECK} & Mini-toets op het einde van elk level. \\
\hline
\end{tabular}

\begin{letop}
\textbf{Schrijf je werkwijze altijd op.} Een antwoord zonder berekening of uitleg is niet volledig. Eindig bij een context steeds met een \textbf{antwoordzin}: zo toon je dat je het getal ook begrijpt.
\end{letop}

\pagina

\subsection*{Doelenkaart: wat moet je kunnen?}

Vul de kaart in \textbf{voor} je begint en opnieuw \textbf{na} de eindtoets. Gebruik: \textbf{+} (lukt zeker), \textbf{±} (lukt half), \textbf{–} (lukt nog niet).

\breedtes{7/69/9/9}
\begin{tabular}{llll}
\hline
LPD & Ik kan… & Voor & Na \\
\hline
 & de begrippen populatie, steekproef, variabele en eenheid gebruiken, het soort data bepalen en centrum- en spreidingsmaten berekenen en interpreteren (springplank, level 0 – 2). & & \\
11 & verklaren waarom representativiteit belangrijk is bij een steekproef, en vertekende steekproeven (selectiebias) herkennen. & & \\
12 & in concrete situaties het verschil uitleggen tussen samenhang en causaliteit, en observationeel van experimenteel onderzoek onderscheiden. & & \\
13 & de normale verdeling gebruiken als model bij data: de toepasbaarheid grafisch beoordelen, $\mu$ en $\sigma$ schatten met $\overline{x}$ en $s$ en de betekenis van $\mu$ en $\sigma$ in de Gausskromme uitleggen. & & \\
14 & kansen berekenen bij een normaal verdeelde variabele. & & \\
1 – 4 & fenomenen uit de realiteit beschrijven met wiskundige begrippen (1), problemen oplossen met een goede aanpak en een antwoordzin (2), ICT gebruiken voor berekeningen en grafieken (3) en mijn redenering uitleggen (4). & & \\
\hline
\end{tabular}
